\documentclass[12pt]{extarticle}
\usepackage{array}
\usepackage[T1]{fontenc}
\usepackage[Russian]{babel}
\usepackage{booktabs}
\usepackage{graphicx}
\usepackage{spreadtab}
\usepackage{numprint}
\usepackage[margin=3cm]{geometry}
\usepackage{float}
\usepackage{xfp}
\usepackage{amsmath} 
\usepackage{indentfirst}  
\begin{document}
\author{Осипов Александр Б06-604}
\title{Лабораторная работа 1.3.1. Определение модуля Юнга на основе исследования деформаций растяжения и изгиба}
\maketitle

\section{Цель}

Экспериментально получить зависимость между напряжением и деформацией (закон Гука) для одноосного растяжения и изгиба.
\vspace{5mm}

\section{Оборудование}

Первая часть: прибор Лермантова, проволки из исследуемого материала, зрительная труба со шкалой, набор грузов.

Вторая часть: стойка для изгибания балки, индикатор для измерения величины прогиба.

\vspace{5mm}

\section{Теоретическая часть (приницип работаы счётчика)}

\vspace{5mm}

\section{Ход работы}

\subsection{Часть 1}

Измерение диаметра проволки
\begin{spreadtab}{{tabular}{ccccccccccc}}
    1&2&3&4&5&6&7&8&9&10&@\text{mean}\\
    1&2&3&4&5&6&7&8&9&10&@\text{mean}\\
\end{spreadtab}

\vspace{5mm}
Длина проволки:
l = 

Расстояние от шкалы до зеркальца:
h = 

\vspace{5mm}
Зависимость деформации от нагрузки:
\begin{tabular}{cccccccccc}
    1&2&3&4&5&6&7&8&9&10
\end{tabular} 

\vspace{5mm}
Определение материала проволки:


\subsection{Часть 2}

Расстояние между рёбрами призм:
w = 50,4 см

Балка:
Ширина 21,4мм
Толщина 4,4 мм

\vspace{5mm}
Полный оборот прибора - 1 мм
Начал c 7 оборотов 
основная шкала - одно деление 0,01 мм
Деформация балки:

Нагрузка по центру: 


\begin{spreadtab}{{tabular}{ccccccccccc}}
    @номер замера & @масса груза, г делить на 10& @смещение стрелы прибора (деления) & @\text{mean}\\
    1& 5087& 67 & \\
    2& 4618& 60 & \\
    3& 5000& 69 & \\
    4& 5000+5087 & 140 & \\
    5& 5000+5087+5030 & 208 & \\
    6& 5000+5087+5030+4679+5035 & 340 & \\
    7& 5000+5087+5030+4679+5035+4618+4636 & 473 & \\
\end{spreadtab}


нагрузка 27 слева см
\begin{spreadtab}{{tabular}{ccccccccccc}}
    @номер замера & @масса груза, г  делить на 10& @смещение стрелы прибора (деления) & @\text{mean}\\
    1& 5035& 69 & \\
    2& 5035+4679& 130 & \\
    3& 5035+4679+4618& 198 & \\
    4& 5035+4679+4618+5087 & 261 & \\
    5& 5035+4679+4618+5087+4723+5020 & 403 & \\
    6& 5035+4679+4618+5087+4723+5020+5030 & 466 & \\
\end{spreadtab}


Деревяная балка 1
ширина 20мм
толщина 10,1
В среднем балка при убирании груза невозвращалась на 3-4 деления до нуля (0,03-0,04 мм)
Деформация балки (после переворота):


Деревяная балка 1
ширина 20мм
толщина 10,1

по центру
\begin{spreadtab}{{tabular}{ccccccccccc}}
    @номер замера & @масса груза, г делить на 10 & @смещение стрелы прибора (деления) & @\text{mean}\\
    1& 5035& 39 & \\
    2& 5035+4679& 87 & \\
    3& 5035+4679+4618& 138 & \\
    4& 5035+4679+4618+5087 & 189 & \\
    5& 5035+4679+4618+5087+4723+5020 & 284 & \\
    6& 5035+4679+4618+5087+4723+5020+5030 & 339 & \\
\end{spreadtab}


30 см слева
\begin{spreadtab}{{tabular}{ccccccccccc}}
    @номер замера & @масса груза, г делить на 10 & @смещение стрелы прибора (деления) & @\text{mean}\\
    1& 5035& 34 & \\
    2& 5035+4679&  & \\
    3& 5035+4679+4618& 128 & \\
    4& 5035+4679+4618+5087 & 175 & \\
    5& 5035+4679+4618+5087+4723+5020 & & \\
    6& 5035+4679+4618+5087+4723+5020+5030 & 318 & \\
\end{spreadtab}

Переворот
\begin{spreadtab}{{tabular}{ccccccccccc}}
    @номер замера & @масса груза, г делить на 10 & @смещение стрелы прибора (деления) & @\text{mean}\\
    1& 5035& 49 & \\
    2& 5035+4679& 98 & \\
    3& 5035+4679+4618&  & \\
    4& 5035+4679+4618+5087 & 176 & \\
    5& 5035+4679+4618+5087+4723+5020 &  & \\
    6& 5035+4679+4618+5087+4723+5020+5030 & 322 & \\
\end{spreadtab}


Деревяная балка 2 
20 мм ширина
9,5 мм толщина

\begin{spreadtab}{{tabular}{ccccccccccc}}
    @номер замера & @масса груза, г делить на 10 & @смещение стрелы прибора (деления) & @\text{mean}\\
    1& 5035& 50\\
    2& 5035+4679&  109& \\
    3& 5035+4679+4618&  & \\
    4& 5035+4679+4618+5087 & 176 & \\
    5& 5035+4679+4618+5087+4723+5020 &  & \\
    6& 5035+4679+4618+5087+4723+5020+5030 & 331 & \\
\end{spreadtab}




ЧАсть 2 Эксперимент с проволкой 

d(pr)= 0.46 мм

рычаг прибора r = 15 мм

расстояние 124 см от шкалы до зеркальца

177 см длина проволки

Эксперимент 1 нагрузка увеличивается
\begin{spreadtab}{{tabular}{ccccccccccc}}
    @номер замера & @масса груза, г делить на 10 & @смещение по шкале в см (деления) & @\text{mean}\\
    1& 5044& 19,6 &\\
    2& 5044+2456&  22,3& \\
    3& 5044+2456+2453+2455& 27,3 & \\
    4& 5044+2456+2453+2455+2461+2458 & 31,8 & \\
    5& 5044+2456+2453+2455+2461+2458+2453+2456 & 36,2 & \\
    6& 5044+2456+2453+2455+2461+2458+2453+2456+2455+2456 & 40,6 & \\
    7& 5044+2456+2453+2455+2461+2458+2453+2456+2455+2456+2456 & 42,8 & \\
\end{spreadtab}


Эксперимент 2 нагрузка уменьшается (читать снизу вверх)
\begin{spreadtab}{{tabular}{ccccccccccc}}
    @номер замера & @масса груза, г делить на 10 & @смещение по шкале в см (деления) & @\text{mean}\\
    1& 5044& 19,5 &\\
    2& 5044+2456&  22,5& \\
    3& 5044+2456+2453+2455& 27,2 & \\
    4& 5044+2456+2453+2455+2461+2458 & 31,8 & \\
    5& 5044+2456+2453+2455+2461+2458+2453+2456 & 36,3 & \\
    6& 5044+2456+2453+2455+2461+2458+2453+2456+2455+2456 & 40,7 & \\
    7& 5044+2456+2453+2455+2461+2458+2453+2456+2455+2456+2456 & 42,8 & \\
\end{spreadtab}


Эксперимент 1.3.2(2)
377,5 и 373,1 г это грузы

1750+-2 мм длина
3,01 мм
Эксперимент 2 нагрузка уменьшается (читать снизу вверх)
\begin{spreadtab}{{tabular}{ccccccccccc}}
    @номер замера & @длина до груза, см & @амплитуда в ед на шкале &@количество периодов (время 10 колебаний делить на 10)& @\text{mean}\\
    1 & 3 & 10 & 2,384 & \\
    2 & 3 & 5 & 2,349 & \\
    3 & 5 & 2,883 & \\
    4 & 7 & 3,392 & \\
    5 & 9 & 3,984 & \\
\end{spreadtab}

\end{document}
